% Final assembly of Kummer's criterion.

This chapter combines the minus criterion of
Chapter~\ref{ch:hminus-formula} with the cyclotomic-unit results of
Chapter~\ref{ch:cyclotomic-units}.  The bridge is the implication
$p \mid h^+ \Rightarrow p \mid h^-$: if $p \nmid h^-$, the minus criterion
clears every Bernoulli numerator in Kummer's range, the Kummer logarithm
determinant is then nonzero, $C^+$ is $p$-saturated, and the
prime-conductor index theorem forces $p \nmid h^+$.

\section{\texorpdfstring{The implication from $h^+$ to $h^-$}{The implication from h-plus to h-minus}}

\begin{theorem}\label{thm:bernoulli-nonzero-index}
  \lean{KummerCriterion.not_dvd_cyclotomicUnitIndex_of_bernoulli_nonzero}\leanok
  \uses{thm:kummer-log-det-bernoulli, thm:kummer-log-det-saturation,
  thm:not-dvd-index-of-pSaturated}
Let $p$ be an odd prime.  Suppose that
\[
  \forall j,\ 1\le j,\ 2j\le p-3,\quad
  p\nmid\operatorname{num}(B_{2j}).
\]
Then
\[
  p\nmid[(\mathcal O_{K^+})^\times:C^+].
\]
\end{theorem}
\begin{proof}\leanok
If $p=3$, Kummer's range is empty, $K^+ = \Q$, the group $C^+$ is the full
unit group $\{\pm1\}$, and the index is $1$.  For $p\ge5$,
\cref{thm:kummer-log-det-bernoulli} identifies the Bernoulli nonvanishing
hypothesis with the nonvanishing of the Kummer logarithm determinant.  By
\cref{thm:kummer-log-det-saturation}, the subgroup $C^+$ is then
$p$-saturated in the full unit group, and
\cref{thm:not-dvd-index-of-pSaturated} concludes that the index is prime
to $p$.
\end{proof}

\begin{theorem}\label{thm:not-dvd-hplus-of-not-dvd-hminus-units}
  \lean{KummerCriterion.not_dvd_hPlus_of_not_dvd_hMinus_units}\leanok
  \uses{thm:hminus-nondiv-to-bernoulli-nondiv, thm:bernoulli-nonzero-index,
  thm:sinnott-index-pprimary, thm:CPlus-normalized-index}
If $p\nmid h^-(K)$, then $p\nmid h^+(K)$.
\end{theorem}
\begin{proof}\leanok
By \cref{thm:hminus-nondiv-to-bernoulli-nondiv}, the hypothesis
$p\nmid h^-(K)$ implies that no Bernoulli numerator in Kummer's range is
divisible by $p$, so \cref{thm:bernoulli-nonzero-index} gives
\[
  p\nmid[(\mathcal O_{K^+})^\times:C^+].
\]
Suppose for contradiction that $p\mid h^+(K)$.  The prime-conductor index
theorem (\cref{thm:sinnott-index-pprimary}) then gives $p \mid
[(\mathcal O_{K^+})^\times:C^+_{\mathrm{norm}}]$, and by
\cref{thm:CPlus-normalized-index} this is equivalent to $p$ dividing the
ordinary $C^+$ index --- a contradiction.
\end{proof}

\begin{theorem}\label{thm:hplus-to-hminus-units}
  \lean{KummerCriterion.weakReflection_dvd_hMinus_of_dvd_hPlus_units}\leanok
  \uses{thm:not-dvd-hplus-of-not-dvd-hminus-units}
If $p\mid h^+(K)$, then $p\mid h^-(K)$.
\end{theorem}
\begin{proof}\leanok
Contrapositive of \cref{thm:not-dvd-hplus-of-not-dvd-hminus-units}: if
$p\nmid h^-(K)$, that theorem gives $p\nmid h^+(K)$.
\end{proof}

\section{The total class-number criterion}

\begin{theorem}\label{thm:dvd-h-iff-bernoulli-units}
  \lean{KummerCriterion.dvd_h_iff_exists_dvd_bernoulli_units}\leanok
  \uses{thm:dvd-h-iff-dvd-hminus-from-plus-to-minus,
  thm:hplus-to-hminus-units, thm:hminus-bernoulli-final}
For an odd prime $p$ and the $p$th cyclotomic field $K$,
\[
  p\mid h(K)
  \quad\Longleftrightarrow\quad
  \exists k,\ 1\le k,\ 2k\le p-3,\quad
  p\mid\operatorname{num}(B_{2k}).
\]
\end{theorem}
\begin{proof}\leanok
The implication $p\mid h^+(K)\Rightarrow p\mid h^-(K)$ is
\cref{thm:hplus-to-hminus-units}; feeding it to
\cref{thm:dvd-h-iff-dvd-hminus-from-plus-to-minus} converts $p\mid h(K)$
into $p\mid h^-(K)$.  The minus class-number criterion
(\cref{thm:hminus-bernoulli-final}) rewrites the latter as the existence of
a Bernoulli numerator in Kummer's range divisible by $p$.
\end{proof}

\section{Kummer's criterion}

\begin{theorem}\label{thm:kummer-criterion}
  \lean{KummerCriterion}\leanok
  \uses{def:regular-prime, thm:dvd-h-iff-bernoulli-units}
Let $p$ be an odd prime.  Then
\[
  p\ \text{is regular}
  \quad\Longleftrightarrow\quad
  \forall k,\ 1\le k,\ 2k\le p-3,\quad
  p\nmid\operatorname{num}(B_{2k}).
\]
\end{theorem}
\begin{proof}\leanok
By \cref{def:regular-prime}, $p$ is regular if and only if $p$ is coprime
to the class number of the $p$th cyclotomic field; since $p$ is prime,
this is equivalent to $p\nmid h(K)$.  By
\cref{thm:dvd-h-iff-bernoulli-units}, the negation of $p\mid h(K)$ is
exactly the assertion that no $k$ in the range $1\le k$, $2k\le p-3$
satisfies $p\mid\operatorname{num}(B_{2k})$, which is the displayed
universal nondivisibility.
\end{proof}
